\section*{GAP code}

\subsection*{Multiplication tables}

This is a naive GAP code that produces the table
shown on page~\pageref{multiplication_table}:
\begin{lstlisting}
gap> s := [0,1,2,3];;
gap> m := NullMat(4,4);;
gap> for x in s do
> for y in s do
> m[x+1][y+1] := (x+y) mod 4;
> od;
> od;
gap> Display(m);
[ [  0,  1,  2,  3 ],
  [  1,  2,  3,  0 ],
  [  2,  3,  0,  1 ],
  [  3,  0,  1,  2 ] ]
\end{lstlisting}


\subsection*{Permutations in GAP}

In our calculations with permutations, we have adopted the right-to-left
convention, aligning with the rest of the course, where we consider groups
acting on the left on sets (see~\pageref{convention:right-to-left}).  However,
some people prefer the left-to-right convention, which corresponds to the way
English is read (and to groups acting on the right). This is the convention
used by GAP:
\begin{lstlisting}
gap> S3 := SymmetricGroup(3);;
gap> a := (1,2);;
gap> b := (2,3);;
gap> a*b;
(1,3,2)
\end{lstlisting}
But there's no need to worry: both conventions are equivalent, see Example \ref{exa:op}.

\subsection*{Example~\ref{exa:dihedral}}

In this example, we show a matrix representation
for dihedral groups. To help to fix ideas,
we also show what happens with the dihedral group
$\D_4$ of eight elements. Here is the GAP
code to construct this particular group.
\begin{lstlisting}
gap> r := [[0,-1],[1,0]];;
gap> s := [[1,0],[0,-1]];;
gap> G := Group(r,s);;
gap> Order(G);
8
gap> r^3*s;
[ [ 0, -1 ], [ -1, 0 ] ]
\end{lstlisting}


\subsection*{Exercise~\ref{xca:D3}}

The following GAP code uses some more advanced methods to construct elements such as $\sqrt{3}$. These details are not needed for solving the exercise and may be ignored; they are included only to record a few useful tips for certain constructions in GAP.

\begin{lstlisting}
gap> z := Sqrt(3);;
gap> z^2;
3
gap> r := [[-1/2,-z/2],[z/2,-1/2]];;
gap> s := [[1,0],[0,-1]];;
gap> G := Group(r,s);;
gap> Order(G);
6
\end{lstlisting}

\subsection*{Example~\ref{exa:Q8}}

This is the code to create the quaternion group $Q_8$
in GAP:
\begin{lstlisting}
gap> i := E(4);;
gap> a := [[i,0],[0,-i]];;
gap> b := [[0,1],[-1,0]];;
gap> G := Group(a,b);;
gap> StructureDescription(G);
"Q8"
\end{lstlisting}

\subsection*{Example~\ref{exa:U(Z/8)again}}

This simple example can be carried out entirely by hand,
but the GAP code includes some tricks that may be useful
later, especially for identifying elements of the group
of units of $\Z/8$ inside $\Z/8$ itself:
\begin{lstlisting}
gap> R := Integers mod 8;;
gap> G := Units(R);;
gap> Set(G, Int);
[ 1, 3, 5, 7 ]
gap> H := Subgroup(G, [ZmodnZObj(3,8)]);;
gap> Set(H, Int);
[ 1, 3 ]
\end{lstlisting}

\subsection*{Example~\ref{exa:Carmichael}}
\label{gap:Carmichael}

A nice subgroup of $\Sym_{16}$
in which the set of commutators is not a subgroup:

\begin{lstlisting}
gap> S16 := SymmetricGroup(16);;
gap> a := (1,3)(2,4);;
gap> b := (5,7)(6,8);;
gap> c := (9,11)(10,12);;
gap> d := (13,15)(14,16);;
gap> e := (1,3)(5,7)(9,11);;
gap> f := (1,2)(3,4)(13,15);;
gap> g := (5,6)(7,8)(13,14)(15,16);;
gap> h := (9,10)(11,12);;
gap> G := Group(a,b,c,d,e,f,g,h);;
gap> D := DerivedSubgroup(G);;
gap> Size(D);
16
gap> comms := Set(Cartesian(G,G), p -> p[1]*p[2]*p[1]^-1*p[2]^-1);;
gap> Size(comms);
15
gap> c*d in Difference(Set(D), comms);
true
\end{lstlisting}

\subsection*{Example~\ref{exa:Guralnick}}
\label{gap:Guralnick}

In the example, we claimed that
this is an example of a group where the set of commutators is not a subgroup. We
start creating a function that returns the set of commutators:
\begin{lstlisting}
gap> SetOfCommutators := function(G)
> return Set(Cartesian(G,G), p -> Comm(p[1],p[2]));
> end;;
\end{lstlisting}
Now we check that in Guralnick's group, the commutator subgroup
differs from the set of commutators:
\begin{lstlisting}
gap> S12 := SymmetricGroup(12);;
gap> a := (1,3,5)(2,4,6)(7,11,9)(8,12,10);;
gap> b := (3,9,4,10)(5,8)(6,7)(11,12);;
gap> G := Group(a,b);;
gap> Order(G);
96
gap> D := DerivedSubgroup(G);;
gap> Order(D);
32
gap> Size(SetOfCommutators(G));
29
\end{lstlisting}
Let us see that no smaller group
have the desired property:
\begin{lstlisting}
gap> allsmall := Concatenation(List([1..95], AllSmallGroups));;
gap> Filtered(allsmall, G -> Size(DerivedSubgroup(G))
> <> Size(SetOfCommutators(G)));
[  ]
\end{lstlisting}



\subsection*{Exercise~\ref{xca:orders_Z6}}

Here is a GAP solution, where for example we see that, as we know,
$1\in\Z/6$ is a generator,
$2\in\Z/6$ has order three, and $4\in\Z/6$ has order three:
\begin{lstlisting}
gap> G := CyclicGroup(IsPermGroup, 6);;
gap> gen := GeneratorsOfGroup(G)[1];;
gap> Set([0..5], k -> [k, Order(gen^k)]);
[ [ 0, 1 ], [ 1, 6 ], [ 2, 3 ], [ 3, 2 ], [ 4, 3 ], [ 5, 6 ] ]
\end{lstlisting}

\subsection*{Exercise~\ref{xca:order_ab}}

Let us see a GAP solution:
\begin{lstlisting}
gap> a := [[1,-1],[1,0]];;
gap> b := [[0,1],[-1,-1]];;
gap> Order(a);
6
gap> Order(b);
3
gap> IsFinite(Group(a*b));
false
\end{lstlisting}

\subsection*{Example~\ref{exa:S3_cosets}}

As we discussed, the annoying fact that GAP multiplies permutations in the opposite order from us is not a serious issue. In this particular example, we only need to work with right cosets instead of left cosets.
Why? Because switching from left to right cosets exactly compensates for the reversed convention in permutation multiplication!

\begin{lstlisting}
gap> G := SymmetricGroup(3);;
gap> H := Subgroup(G,[(1,2)]);;
gap> K := Subgroup(G,[(1,2,3)]);;
gap> Set(H, h -> h*(1,2,3));
[ (1,2,3), (1,3) ]
gap> Set(H, h -> h*(2,3));
[ (2,3), (1,3,2) ]
\end{lstlisting}

\subsection*{Bonus exercise~\ref{bonus:SL23}}

With GAP the exercise becomes almost trivial, which is why I encourage the reader to try solving it first without using a computer.
\begin{lstlisting}
gap> G := SL(2,3);;
gap> Size(Filtered(AllSubgroups(G), H -> Size(H) = 12));
0
\end{lstlisting}

\subsection*{Example~\ref{exa:normal_A4}}

For GAP, getting the list of normal subgroups of $\Alt_4$
is easy:
\begin{lstlisting}
gap> G := AlternatingGroup(4);;
gap> NormalSubgroups(G);
[ Alt( [ 1 .. 4 ] ), Group([ (1,4)(2,3), (1,2)(3,4) ]), Group(()) ]
\end{lstlisting}


\subsection*{Exercise~\ref{xca:D4_normal}}

The exercise shows that normality is in general not transitive:
\begin{lstlisting}
gap> G := DihedralGroup(8);;
gap> gens := GeneratorsOfGroup(G);;
gap> s := gens[1];; r := gens[2];;
gap> Order(r);
4
gap> Order(s);
2
gap> N := Subgroup(G,[s,r^2]);;
gap> H := Subgroup(G,[s]);;
gap> IsNormal(N,H);
true
gap> IsNormal(G,N);
true
gap> IsNormal(G,H);
false
\end{lstlisting}

\subsection*{Example~\ref{exa:HKneKH}}

Here is the GAP code:
\begin{lstlisting}
gap> G := SymmetricGroup(3);;
gap> H := Subgroup(G,[(1,2)]);;
gap> K := Subgroup(G,[(2,3)]);;
gap> Set(Cartesian(H,K), p -> p[1]*p[2])
> = Set(Cartesian(K,H), p -> p[1]*p[2]);
false
\end{lstlisting}

\subsection*{Example~\ref{exa:anotherHK}}

Here is the code to verify some of the statements mentioned:
\begin{lstlisting}
gap> G := SymmetricGroup(3);;
gap> H := Subgroup(G,[(1,2)]);;
gap> K := Subgroup(G,[(1,2,3)]);;
gap> IsNormal(G,K);
true
gap> Size(Set(Cartesian(H,K), p -> p[1]*p[2]));
6
gap> Size(Intersection(H,K));
1
\end{lstlisting}

\subsection*{Example~\ref{exa:U(21)}}

Recall that we work multiplicatively with the group of units. Here is the code:
\begin{lstlisting}
gap> R := Integers mod 21;;
gap> G := Units(R);;
gap> Set(G, Int);
[ 1, 2, 4, 5, 8, 10, 11, 13, 16, 17, 19, 20 ]
gap> f := GroupHomomorphismByFunction(G, G, x -> x^3);;
gap> Set(Kernel(f), Int);
[ 1, 4, 16 ]
gap> Set(Image(f), Int);
[ 1, 8, 13, 20 ]
\end{lstlisting}
Alternative code:
\begin{lstlisting}
gap> SetInfoLevel(InfoWarning, 0);;
gap> U := Filtered([0..20], x -> Gcd(x,21) = 1);
[ 1, 2, 4, 5, 8, 10, 11, 13, 16, 17, 19, 20 ]
gap> G := Group(List(U, x -> ZmodnZObj(x,21)));;
gap> f := GroupHomomorphismByFunction(G, G, x -> x^3);;
gap> Set(Image(f), Int);
[ 1, 8, 13, 20 ]
\end{lstlisting}

% How to check that a map is a homomorphism of abelian groups?
% gap> P := Image(IsomorphismPermGroup(G));
% gap> f := GroupHomomorphismByImages(P,P,GeneratorsOfGroup(P),
% > List(GeneratorsOfGroup(P), x -> x^3));
% gap> IsGroupHomomorphism(f);
% true

\subsection*{Example~\ref{exa:order4}}

The groups are not isomorphic because
they do not have the same number of elements of order four:
\begin{lstlisting}
gap> A := AbelianGroup([4]);;
gap> B := AbelianGroup([2,2]);;
gap> Size(Filtered(A, x -> Order(x) = 4));
2
gap> Size(Filtered(B, x -> Order(x) = 4));
0
\end{lstlisting}

\subsection*{Example~\ref{exa:U(Z/5)}}

Both groups are cyclic of order four:
\begin{lstlisting}
gap> A := Units(Integers mod 5);;
gap> B := Units(Integers mod 10);;
gap> IsCyclic(A);
true
gap> IsCyclic(B);
true
gap> Size(A);
4
gap> Size(B);
4
gap> IsomorphismGroups(A,B) <> fail;
true
\end{lstlisting}

\subsection*{Exercise~\ref{xca:U(Z/10)}}

The groups are not isomorphic. Why?
Here is the code:
\begin{lstlisting}
gap> A := Units(Integers mod 10);;
gap> B := Units(Integers mod 12);;
gap> IsomorphismGroups(A,B) = fail;
true
gap> StructureDescription(A);
"C4"
gap> StructureDescription(B);
"C2 x C2"
\end{lstlisting}

\subsection*{Example~\ref{exa:Z24}}

We first define the group of integers modulo 24,
and then use a map to represent the elements
of our group correctly inside $\Z/24$:
\begin{lstlisting}
gap> n := 24;;
gap> G := CyclicGroup(IsPermGroup, n);;
gap> gen := GeneratorsOfGroup(G)[1];;
gap> e := function(k) return gen^(k mod n); end;;
gap> val := function(x) 
> return First([0..n-1], k -> gen^k = x); 
> end;;
gap> H := Subgroup(G, [e(4)]);;
gap> Set(H, val);
[ 0, 4, 8, 12, 16, 20 ]
gap> N := Subgroup(G, [e(6)]);;
\end{lstlisting}
Now we check the claims on subgroups
related to $H$ and $N$:
\begin{lstlisting}
gap> H = Subgroup(G,[e(4)]);
true
gap> N = Subgroup(G,[e(6)]);
true
gap> ClosureGroup(H,N) = Subgroup(G,[e(2)]);
true
gap> Set(Intersection(H,N), val);
[ 0, 12 ]
\end{lstlisting}
In the example there are also some claims about (left)
cosets of $N$ and $H\cap N$:
\begin{lstlisting}
gap> Set(N, x -> val(e(0)*x));
[ 0, 6, 12, 18 ]
gap> Set(N, x -> val(e(2)*x));
[ 2, 8, 14, 20 ]
gap> Set(N, x -> val(e(4)*x));
[ 4, 10, 16, 22 ]
gap> HmeetN := Intersection(H,N);;
gap> Set(HmeetN, x -> val(e(0)*x));
[ 0, 12 ]
gap> Set(HmeetN, x -> val(e(4)*x));
[ 4, 16 ]
gap> Set(HmeetN, x -> val(e(8)*x));
[ 8, 20 ]
\end{lstlisting}

% > J := H+N;
% > IsIsomorphic(J/sub<J|N>, H/sub<H|H meet N>);
% true

\subsection*{Example~\ref{exa:Q8normal}}

Checking in GAP that every subgroup of $Q_8$
is normal is very easy. Here is the code:
\begin{lstlisting}
gap> G := QuaternionGroup(8);;
gap> ForAll(AllSubgroups(G), S -> IsNormal(G,S));
true
\end{lstlisting}

\subsection*{Example~\ref{exa:Z12->Z6}}

Let us use GAP to illustrate the correspondence
theorem. Instead of using
additive abelian groups we will
use their corresponding permutation groups, as
several algorithms do not work
on abelian groups.
\begin{lstlisting}
gap> C12 := CyclicGroup(IsPermGroup, 12);;
gap> C6 := CyclicGroup(IsPermGroup, 6);;
gap> f := GroupHomomorphismByImages(C12, C6,
> GeneratorsOfGroup(C12), GeneratorsOfGroup(C6));;
gap> K := Kernel(f);;
gap> IsGroupHomomorphism(f);
true
gap> a := function(x)
> return First([0..11], k -> GeneratorsOfGroup(C12)[1]^k = x);
> end;;
gap> for x in AllSubgroups(C12) do
> if IsSubset(x, K) then
> Set(x, a);
> fi;
> od;
\end{lstlisting}

\subsection*{Exercise~\ref{xca:aut(S3)}}

Again, this is quite straightforward if we are allowed to use GAP.
Here is one possible solution:
\begin{lstlisting}
gap> S3 := SymmetricGroup(3);;
gap> aut_S3 := AutomorphismGroup(S3);;
gap> StructureDescription(aut_S3);
"S3"
\end{lstlisting}
The function \lstinline{StructureDescription} may fail for larger groups,
so it is useful to provide an alternative solution using \lstinline{IsomorphismGroups}.
However, this function requires permutation groups as input,
so we must first convert our automorphism group:
\begin{lstlisting}
gap> P := Image(IsomorphismPermGroup(aut_S3));;
gap> bool := IsomorphismGroups(P,S3) <> fail;;
gap> bool;
true
\end{lstlisting}

\subsection*{Example~\ref{exa:T}}

Let us construct this semidirect product in GAP.
We build our groups as permutation groups, which is necessary
for the type of computations we will perform.
\begin{lstlisting}
gap> Q := CyclicGroup(IsPermGroup, 4);;
gap> K := CyclicGroup(3);;
gap> autK := AutomorphismGroup(K);;
gap> rho := GeneratorsOfGroup(autK)[1];;
gap> IsOne(rho);
false
gap> tau := GroupHomomorphismByImages(Q, autK,
> GeneratorsOfGroup(Q), [rho]);;
gap> T := SemidirectProduct(Q, tau, K);;
gap> Size(T);
12
\end{lstlisting}
Let us quickly check that this group is not isomorphic to $\Alt_4$:
\begin{lstlisting}
gap> A4 := AlternatingGroup(4);;
gap> IsomorphismGroups(T,A4) = fail;
true
\end{lstlisting}

\subsection*{Example~\ref{exa:order21}}

Let us check with GAP that
\[
\langle x,y:x^7=y^3-1,yx=x^2y\rangle\simeq\left\langle
\begin{pmatrix}
    1&1\\
    0&1
\end{pmatrix},
\begin{pmatrix}
    2&0\\
    0&1
\end{pmatrix}\right\rangle.
\]
There are several ways to verify this. In our situation, since we will construct
the groups in different ways, we will simply check
that the outputs of the function \lstinline{IdGroup}
are the same for both groups:
\begin{lstlisting}
gap> F := FreeGroup("x","y");;
gap> x := F.1;; y := F.2;;
gap> G := F / [x^7, y^3, y*x/(x^2*y)];;
gap> one := Z(7)^0;; zero := 0*Z(7);; two := 2*one;;
gap> H := Group([[one,one],[zero,one]], [[two,zero],[zero,one]]);;
gap> IdGroup(G) = IdGroup(H);
true
\end{lstlisting}


\subsection*{Examples~\ref{exa:[6,100,45]}}

Computing abelian invariants of abelian groups is quite easy:
\begin{lstlisting}
gap> G := AbelianGroup([6,100,45]);;
gap> Filtered(ElementaryDivisorsMat(DiagonalMat([6,100,45])),
> d -> d > 1);
[ 30, 900 ]
\end{lstlisting}

In the same way, we can easily verify the abelian invariants computed in
Example~\ref{exa:[10,15,20,25]} and get the
solution to Exercise~\ref{xca:factors:24,12,4,2}.

\subsection*{Exercise~\ref{xca:howmany:500}}

One way to obtain the solution is to inspect the groups of order 500 in the
\lstinline{SmallGroups} library:
\begin{lstlisting}
gap> Number(AllSmallGroups(200), IsAbelian);
6
\end{lstlisting}

% \begin{lstlisting}
% gap> NumberOfSylowSubgroups := function(G,p)
% >   local P;
% >   P := SylowSubgroup(G,p);
% >   return Index(G,Normalizer(G,P));
% > end;;
% gap> Set(Filtered(AllSmallGroups(12),
% > G -> NumberOfSylowSubgroups(G,3) <> 1), StructureDescription);
% [ "A4" ]
% \end{lstlisting}
